\chapter{What a Symbol System Would Have to Be}
\label{ch:notation}
% ===================================================================

Chapter~\ref{ch:trick} showed the method and named the price.
This chapter is the price list.

It is the most technical chapter in the Prestige, and it is
technical on purpose.
The Turn spent three hundred pages establishing that a Kantian
conclusion with a ledger attached is a claim, and that the same
conclusion without one is a posture.
Having just conceded that the book's grounding was analytic rather
than earned, the least i can do is be specific about what earning
it would involve.

The chapter has four movements.
What a symbol system \emph{is}, which turns out to be a prior
question nobody in the grounding literature answers
(Section~\ref{sec:goodman}).
Where the machinery already in hand meets that account and where it
falls short (Section~\ref{sec:bisim}).
What a budget buys, which turns out to be the deepest of three
instances of one move that Chapter~\ref{ch:kant} will gather
(Section~\ref{sec:budget}).
And a construction --- $\RHO[\Bee]$, the calculus-making functor
--- which organises several of the conditions attractively, lands
on a category the Turn already built, and then fails, for a reason
i will make as sharp as i can (Sections~\ref{sec:rho-builtins}
to~\ref{sec:notasolution}).

\section{The prior question}
\label{sec:goodman}

The reference bootstrap problem asks us to construct a symbol
system.
It is therefore prior to ask what one is.

The literature is remarkably thin on this.
Symbol grounding work generally treats \emph{symbol} as understood
and asks how to attach it to sensors.
Representation learning treats discreteness as an architectural
choice, to be imposed with a quantiser where convenient.
Emergent communication treats a symbol as whatever the channel
happens to carry.
None of these supplies a criterion by which a proposed system could
be judged \emph{inadequate}, which is what one needs if the
construction is to be a construction rather than a hope.

Goodman's \emph{Languages of Art} does supply one
\cite{goodman1968}.
It offers a set of conditions distinguishing notational systems
from other symbol systems, developed for the score / performance
relation in music and dance and generalised from there.
i adopt it provisionally --- provisionally because it has known
failure modes, taken up in Section~\ref{sec:failures}, and because
adopting it is a bet that its conditions are the right ones rather
than merely the first ones available.
The one prior attempt to put the theory to computational work is
Goel's, which derives from notationality a class of physical
notational systems and argues that these are exactly the dynamical
systems capable of realising computational information processing
\cite{goel1991}.
That work predates everything in the current setting and has not
been revisited for learned representations.
So far as i can determine it is the whole of the prior art, which
is itself a small scandal.

\subsection{The conditions}
\label{sec:conditions}

Goodman distinguishes \emph{marks} (physical inscriptions or
utterances), \emph{characters} (classes of marks that count as the
same), and \emph{compliants} (the things a character applies to,
its compliance class).
A score is a character.
A performance is a compliant.

A \emph{notational scheme} satisfies two syntactic conditions.

\begin{condition}[Syntactic disjointness]
\label{cond:sd}
No mark belongs to more than one character.
\end{condition}

\begin{condition}[Syntactic finite differentiation]
\label{cond:sfd}
For any two characters $K \neq K'$ and any mark $m$ not belonging
to both, it is theoretically possible to determine that $m$ does
not belong to $K$, or that it does not belong to $K'$.
\end{condition}

\noindent
A \emph{notational system} is a scheme satisfying three further
semantic conditions.

\begin{condition}[Unambiguity]
\label{cond:ua}
Each character has a single compliance class: not different
compliants at different times or in different contexts.
\end{condition}

\begin{condition}[Semantic disjointness]
\label{cond:semd}
No two characters share a compliant.
\end{condition}

\begin{condition}[Semantic finite differentiation]
\label{cond:semfd}
The analogue of Condition~\ref{cond:sfd} on the compliance side.
\end{condition}

\noindent
The two finite-differentiation conditions are what exclude
\emph{dense} systems.
A system whose characters are ordered so that between any two there
is a third fails Condition~\ref{cond:sfd}, and this is Goodman's
account of the difference between a score and a sketch, or between
the digital and the analogue.

\subsection{The roundtrip is derived, not primitive}
\label{sec:roundtrip}

It is natural --- and it was my own first reading --- to take the
criterion of adequacy to be losslessness of the roundtrip: from
score to performance and back by transcription, and from
performance to score and back, without loss of identity.
That is indeed the property Goodman is after and it is the property
that matters here.
But it is a \emph{consequence} of
Conditions~\ref{cond:sd}--\ref{cond:semfd} rather than an
additional stipulation, and the distinction is load-bearing.
The roundtrip tells you when you have succeeded.
The five conditions tell you what to check, and they decompose the
requirement into parts that turn out to have very different
relationships to the machinery this book has built.

\begin{remark}[What the framework does and does not do]
\label{rem:invariance}
Goodman's conditions are conditions on a symbol \emph{system}: a
scheme together with a compliance relation.
They presuppose the compliance relation as given and adjudicate
whether it is notational.
They therefore cannot \emph{select} a compliance relation, and they
are invariant under relabelling: if $\Sigma$ is a notational system
and $\phi$ a bijection on characters, then $\phi(\Sigma)$ is a
notational system with the same compliance classes differently
named.
Goodman's framework says what a solution to the reference bootstrap
problem would look like.
It cannot by itself produce one.
\end{remark}

\noindent
Remark~\ref{rem:invariance} is why adopting Goodman is half the
battle rather than progress on the bootstrap problem proper.
It is also the formal seed of Observation~\ref{obs:social}: a
framework invariant under relabelling cannot fix reference, and
this one is invariant by construction.

\subsection{Failure modes, and provisional repairs}
\label{sec:failures}

\paragraph{The wrong-note paradox.}
Goodman's compliance relation is exact.
A performance containing a single wrong note is not a performance
of the work.
Since performances without wrong notes are essentially never found,
the theory as stated has no instances in musical practice.
This is the standard objection and it is a real one.
i record the shape of the repair now and return to it in
Section~\ref{sec:budget}: exactness should be relativised to a
resource bound, so that two performances are the same performance
when no affordable observation separates them.
That is not an evasion, because in a cost-accounted setting
``affordable'' is a defined quantity rather than a gesture, and the
move converts an absolute condition into a family of conditions
indexed by budget.

\paragraph{Density is a property of the world, not only of the notation.}
Finite differentiation can fail because the notation is badly
designed, or because the phenomenon being notated is dense and no
notation of it could succeed.
Goodman does not clearly separate these, and for our purposes the
separation is essential: we need to know whether a failure to
notate is a failure of the agent or a fact about the environment.
Question~\ref{q:density} states this as an open problem, and it is
the one on the list i would most like answered.

\paragraph{The ontological commitments are separable.}
Goodman's theory of notation comes attached to strong claims about
the identity of works of art.
Recent work in aesthetics argues that the notational apparatus
survives being detached from those commitments and reinterpreted
epistemically \cite{estetika2026}.
i take that detachment for granted.

\section{Where the machinery meets it, and where it does not}
\label{sec:bisim}

\subsection{The two bisimulations}

The natural translation of the roundtrip requirement into
process-calculus vocabulary is a pair of bisimulations.
Write $\mathcal{S}$ for the space of scores, $\mathcal{P}$ for the
space of performances, $\sem{-} : \mathcal{S} \to
\pop(\mathcal{P})$ for the compliance map taking a score to its
compliance class, and $\tau : \mathcal{P} \to \mathcal{S}$ for
transcription.
The requirement is that
\begin{align}
p \bisim_{\mathcal{P}} p' &\iff \tau(p) \bisim_{\mathcal{S}} \tau(p') \label{eq:rt1}\\
s \bisim_{\mathcal{S}} s' &\iff \sem{s} = \sem{s'} \label{eq:rt2}
\end{align}
with $\bisim_{\mathcal{P}}$ and $\bisim_{\mathcal{S}}$ the
respective behavioural equivalences.
Read one way this says the two kernels agree; read another, that
$\tau$ and $\sem{-}$ are mutually inverse on the quotients.

\subsection{What bisimulation supplies}

Conditions~\ref{cond:sd}, \ref{cond:ua} and~\ref{cond:semd} are
\emph{quotient} conditions and fall out of
\eqref{eq:rt1}--\eqref{eq:rt2} directly.
Syntactic disjointness says the character map is a function on
marks, which is what it means for $\bisim_{\mathcal{S}}$ to be an
equivalence.
Unambiguity says the compliance class of a character does not vary,
which is the statement that $\sem{-}$ is well defined on
$\bisim_{\mathcal{S}}$-classes.
Semantic disjointness says compliance classes partition, which is
\eqref{eq:rt2} read as injectivity on the quotient.

This is a good fit, and it is why the translation is worth making
at all.
Bisimulation is the mathematically mature version of \emph{same
behaviour}, it comes with coinductive proof methods, and in the
reflective setting it is the equivalence the namespace logic
already denotes.

\subsection{Where Goodman adds}
\label{sec:adds}

Conditions~\ref{cond:sfd} and~\ref{cond:semfd} are not quotient
conditions, and bisimulation does not supply them.

\begin{observation}[Separation is not quotient]
\label{obs:sep}
Finite differentiation is a \emph{separation} condition: it
requires not merely that distinct characters be distinct, but that
their distinctness be effectively determinable from a mark.
A bisimulation quotient can be perfectly well defined over a dense
space, in which case every pair of distinct classes is
distinguishable in the limit and no pair is distinguishable by any
finite observation.
Such a system satisfies \eqref{eq:rt1}--\eqref{eq:rt2} and fails
Condition~\ref{cond:sfd}.
\end{observation}

\noindent
That is the precise content of the claim that Goodman is compatible
with the bisimulation approach and adds strictly more.
What he adds is an effectivity requirement that process
equivalence, being defined by a greatest fixed point, is
systematically silent about.

And the point is not academic.
It is exactly where the neural implementations fail, and they fail
there without the vocabulary to say so.

\paragraph{LatPlan and the symbol stability problem.}
LatPlan learns propositional and action symbols from unlabelled
image pairs and plans in the resulting latent space
\cite{asai2018}.
Asai and Kajino subsequently identified what they named the
\emph{symbol stability problem}: the propositions produced by the
state autoencoder behave stochastically, so the same input does not
reliably produce the same proposition \cite{asai2019}.
In Goodman's vocabulary this is a failure of
Condition~\ref{cond:sfd} --- one mark, no effective determination
of which character it belongs to --- discovered empirically, by
people who ran into it as an obstacle to making a learned symbol
system behave like a symbol system.

\paragraph{Vector quantisation as a fix for one condition only.}
The VQ-VAE bottleneck assigns continuous latents to a codebook by
nearest neighbour \cite{oord2017}.
This imposes syntactic disjointness and a separation margin by
construction: read through Goodman, it is an engineering solution
to Conditions~\ref{cond:sd} and~\ref{cond:sfd} and to nothing else.
It says nothing about the compliance side, which is why
quantisation alone does not produce symbols that mean anything.
The subsequent literature's efforts to make the codes ``semantic''
by additional training objectives are attempts to buy
Conditions~\ref{cond:ua}--\ref{cond:semfd} that do not know that is
what they are doing.

\subsection{Why bisimulation metrics are the wrong repair}
\label{sec:metrics}

The reinforcement-learning literature has independently encountered
the rigidity of exact bisimulation and responded with quantitative
relaxations: bisimulation metrics \cite{ferns2004} and their
scalable descendants \cite{gelada2019,zhang2021,castro2021}.
These are the correct response to the wrong-note paradox in that
setting and the wrong response here.

\begin{observation}[The metric smooths what notation needs separated]
\label{obs:metric}
A bisimulation metric replaces the equivalence with a distance and
the quotient with a geometry.
It thereby dissolves precisely the discreteness that
Conditions~\ref{cond:sfd} and~\ref{cond:semfd} require.
A representation shaped by a bisimulation metric is, in Goodman's
classification, closer to a sketch than to a score.
\end{observation}

\noindent
The tension is genuine and not merely terminological.
One cannot have a symbol system and a smooth behavioural geometry
at the same level of description.
What one can have is a discrete system whose \emph{admissible
discriminations} are set by a resource bound.

\section{What a budget buys}
\label{sec:budget}

There is a move this book makes repeatedly and has not yet named.
Something is put out of an agent's reach; one asks whether the
barrier is a matter of logic or of money; and the answer keeps
coming back \emph{money}.
It has appeared twice already.
Hennessy--Milner adequacy characterizes a bisimulation class
positively, by formulae of the agent's own logic, and says nothing
whatever about whether the separating formula can be afforded ---
so what an agent cannot distinguish, it usually cannot distinguish
for want of funds rather than for want of concepts.
And Chapter~\ref{sec:scope-manufacture} shows that a learner does
not merely fail to see the distinctions it cannot afford; it fails
to see that they are there, and reports a world with fewer kinds in
it than the world has.

Here is the third instance, and it is the one that bites hardest,
because it is about which marks are the same mark.
Chapter~\ref{ch:kant} collects all three and argues that what they
amount to is Kant's thing-in-itself surviving translation into
computation as an economic fact rather than a logical one.
i mention that now so the reader can see where this is going; the
argument for it is that chapter's, not mine here.

The mortal-scientist framework already prices observation.
Assays cost.
Refutation is budget-relative.
A verdict may be $\bot$ because the budget ran out rather than
because the world was silent.
That is exactly the ingredient Observation~\ref{obs:sep} says is
missing.

\begin{conjecture}[Finite differentiation as a theorem]
\label{conj:fd}
In a cost-accounted setting with budget $\sigma$, define
$x \bisim^{\sigma} y$ to hold when no assay affordable at $\sigma$
separates $x$ from $y$.
Then the quotient by $\bisim^{\sigma}$ satisfies
Conditions~\ref{cond:sfd} and~\ref{cond:semfd} with the budget as
the separation constant, and the resulting family
$\{\Qs\}_{\sigma}$ is a filtered diagram of notational schemes
refining as $\sigma$ grows.
\end{conjecture}

\noindent
i state this as a conjecture rather than a proposition because two
things are unverified.
First, that $\bisim^{\sigma}$ is transitive: affordability is not
obviously composable, and a chain of pairwise-indistinguishable
states may have distinguishable endpoints, which is the standard
failure of relations of the form ``indistinguishable by a bounded
observer''.
Second, that the cost-accounting monad interacts correctly with the
extension of Section~\ref{sec:rho-builtins}; see
Question~\ref{q:monad}.
If the first fails, the repair is to take the transitive closure
and accept a coarser scheme, at the price of the separation
constant no longer being $\sigma$.

Should Conjecture~\ref{conj:fd} hold, the reading is worth having.
\emph{A poorer agent has a coarser notation.}
Notational adequacy becomes a purchasable property rather than a
binary one, the wrong-note paradox is resolved by making exactness
relative to what can be afforded rather than by abandoning it, and
the whole of Goodman's digital / analogue distinction becomes a
statement about budgets rather than about the metaphysics of
inscription.
That is the same move as the other two, applied one level further
down, to the question of what counts as a symbol at all --- and it
has the same consequence, which is that a wall turns into a
frontier and the frontier has a price on it.

\section{Reflection is score construction, not performance}
\label{sec:rho-builtins}

We now have a criterion and no way of satisfying it.
What follows is an attempt to see how far the machinery of this
book can be pushed toward satisfying it, carried as far as it will
go and then stopped.
It is offered as an analogy.
Section~\ref{sec:notasolution} argues, at some length and against
my own construction, why it is not more than that.

The reflective higher-order calculus takes names to be quoted
processes \cite{meredith2005rho}.
For a process $P$, $\quo{P}$ is a name; for a name $x$, $*x$ is a
process; and $*\quo{P} = P$, $\quo{*x} = x$.
Quotation and dereference are mutually inverse.

It is tempting --- and i made this mistake before correcting it ---
to read that immediately as the score / performance relation, with
Goodman's roundtrip holding by construction.
It is not, and seeing why is what makes the rest of this section
follow rather than intrude.

In pure $\rhoc$, both sides of $\quo{-}$ and $*$ are syntax.
$P$ is a term; $\quo{P}$ is a term-code; $*\quo{P}$ is the same
term again.
Nothing here is a performance.
What the operators do is build scores out of scores: quotation lets
a score be embedded in another score as an atom, and dereference
recovers it.
The construction is \emph{recursive score construction}, and its
losslessness is the statement that the recursion is faithful ---
that no information is lost by descending a level and climbing back
up.
That is a genuine and useful property.
It is also a property internal to the notation.
A notational system all of whose compliants are further
inscriptions is precisely the case Goodman analyses under the
heading of scripts that denote other scripts, and it is not yet a
system that says anything about a world.

\begin{observation}[Two readings of $\quo{-}$ and $*$]
\label{obs:tworeadings}
In $\rhoc = \RHO[\varnothing]$, $\quo{-}$ and $*$ are
score-to-score: they witness the faithfulness of a recursion inside
the notation, and there is no performance in the picture.
The compliance relation is trivial, in the sense that every
character complies with syntax.
\end{observation}

\subsection{Builtins turn the recursion into a correspondence}
\label{sec:turn-of-operators}

Now adjoin a builtin process $b$ drawn from the host world.
The syntax of quotation does not change: $\quo{b}$ is formed
exactly as $\quo{P}$ is.
Its semantics changes completely, because $b$ is not a term of the
calculus.
It is a behaviour in the world the computation is embedded in.

At that point, and only at that point, $\quo{b}$ is a score and $b$
is a performance in Goodman's sense.
$\quo{b}$ is a static, inspectable, transmissible mark that can be
handled by the calculus's own machinery --- sent on channels,
pattern-matched, quantified over in the namespace logic --- while
the thing it codes is an event elsewhere, running under laws the
calculus does not supply.
Dereference $*\quo{b}$ is then not a syntactic identity but a
performance instruction: produce, in the host world, something
indistinguishable from $b$.

\begin{observation}[Where the correspondence enters]
\label{obs:enters}
Adjoining builtins does not add a new operator.
It re-types the existing one.
The same $\quo{-}$ that was recursive score construction over
$\rhoc$-terms becomes, on the builtin sort, a score / performance
correspondence; and the same $*$ that was syntactic recovery
becomes performance.
Goodman's roundtrip requirement therefore attaches to $\rhoc$ at
exactly the point where the calculus is extended by the world, and
nowhere else.
\end{observation}

\noindent
This is why the extension is the interesting object and the pure
calculus is not.
In $\RHO[\varnothing]$ there is nothing to ground and the roundtrip
is a tautology --- which is Observation~\ref{obs:sidestep} again,
arrived at from the side of the notation rather than from the side
of the ontology.

\subsection{Two ways to adjoin the world}
\label{sec:two-ways}

There are two ways to perform the extension, and
Observation~\ref{obs:enters} explains why only one of them does the
job.

\paragraph{Builtin names.}
Adjoin name constants beyond $\quo{\nil}$, understood as channels
to the host world.
This gives an interface but no theory of what is on the other end:
the calculus can send and receive, and has nothing to say about the
behaviour it is interacting with.
It is the applied-$\pi$ move \cite{abadi2001}, where a term algebra
of data is adjoined and the added material is inert with respect to
the operational semantics.
Crucially it does \emph{not} re-type $\quo{-}$: an adjoined name is
a name, not the quotation of anything, so there is no performance
for it to be the score of, and Observation~\ref{obs:tworeadings}
still applies.
The roundtrip stays inside the notation.

\paragraph{Builtin processes.}
Adjoin processes $b$ drawn from the host world.
Then $\quo{b}$ is a symbol \emph{for} $b$, and $*\quo{b}$ is a
process in the host world indistinguishable from $b$.
This is the version that does the work, for the reason just given:
it is the only one of the two under which quotation acquires a
compliance class outside the syntax.

\begin{remark}[This is not psi-calculi]
\label{rem:psi}
The psi-calculi framework parameterises $\pi$-like calculi by
nominal datatypes for terms, conditions and assertions, together
with a channel-equivalence relation, a composition, a unit and an
entailment relation, and subsumes applied $\pi$, spi, fusion and
concurrent constraint $\pi$ as instances \cite{bengtson2011}.
The adjoined material there is \emph{data}: inert under the
operational semantics, transmitted but not run.
Adjoining builtin \emph{processes} is a different move, because the
adjoined objects have their own dynamics and the calculus ceases to
be closed --- part of the reduction relation is delegated to the
host.
Readers who know psi-calculi will otherwise assume this is an
instance of it.
\end{remark}

\section{The calculus-making functor, and a category we already have}
\label{sec:functor}

The natural way to organise this is to stop thinking of $\rhoc$ as
a calculus and start thinking of it as a functor from behaviours to
calculi:
\[
\RHO[\Bee] ::= \nil \mid \Bee \mid \inn{y}{x}{P} \mid \out{x}{Q}
\mid P \Par Q \mid *x,
\qquad x, y ::= \quo{P}
\]
with the usual equations and rewrites, and with $\Bee$'s behaviour
specified in a source category $\Cat$ and imported into a version
of the parallel-composition rule.
The pure calculus is the fibre over the empty behaviour,
$\rhoc = \RHO[\varnothing]$, and everything the Turn did with
$\rhoc$ was done in that fibre.

Now the part i find genuinely satisfying, because it was not
planned.

Chapter~\ref{ch:catgslt} built a category of theories, and built it
carefully, because the obvious definition was wrong three separate
ways.
Objects are GSLTs; morphisms are pseudofunctors of context
multicategories that preserve context-labelled bisimulation, with
the bisimulation on the target computed only over the image of the
source's contexts (Definition~\ref{ob:def:morphism}).
And to stop the resulting comparison being trivial, that chapter
imposed two conditions on a morphism $F : G \to H$:

\begin{itemize}[leftmargin=1.6em]
\item $F$ is \emph{hosting} when the target can run the source ---
the encoding reflects as well as preserves context-labelled
transition structure, so no branching present in $G$ is lost in
$H$;
\item $F$ is \emph{exhausting} when the target is nothing the
source cannot assemble --- every context of $H$ relevant to the
image is, up to equivalence, in the image of a context of $G$.
\end{itemize}

\noindent
Those are faithfulness and density, and they were introduced to
make the expressiveness preorder on calculi mean something.

$\RHO[-]$ is a functor into that category.
Take the morphisms of $\Cat$ to be behaviour-preserving maps.
There is then a unit $\eta_{\Bee} : \Bee \to \RHO[\Bee]$ including
builtins as processes, and:

\begin{itemize}[leftmargin=1.6em]
\item $\eta_{\Bee}$ \emph{hosting} says distinct builtin behaviours
receive distinct symbols and no branching in the world is collapsed
by the calculus, which is performance $\to$ score $\to$ performance
losslessness;
\item $\eta_{\Bee}$ \emph{exhausting} says every calculus-level
distinction over builtins is realised by a builtin distinction,
which is score $\to$ performance $\to$ score losslessness, and
which forbids the calculus from minting symbols the world cannot
cash.
\end{itemize}

\begin{observation}[Goodman's adequacy is hosting and exhausting]
\label{obs:hosting}
Goodman's two-way losslessness for the score / performance relation
of Observation~\ref{obs:enters} is the statement that $\eta_{\Bee}$
is hosting and exhausting in the sense of
Chapter~\ref{ch:catgslt}, with respect to behavioural equivalence.
\end{observation}

\noindent
i want to be careful about how much that is and how much it is not.
It is not a solution to anything.
What it is, is the observation that the pair of conditions this
book introduced to compare one calculus with another turns out to
be the criterion of notational adequacy when one of the two is a
world instead of a calculus.
The question the Turn asked was: how much of calculus $G$ can
calculus $H$ host, and how much of $H$ does $G$ exhaust?
The question here is: how much of world $\Bee$ can $\rhoc$ host,
and does $\rhoc$ mint anything $\Bee$ cannot cash?
Same two conditions, same category, one argument changed from a
theory to a world.
Goodman's criterion of adequacy becomes a property one could try to
prove, which is more than it was before.

\subsection{Three further things the functorial reading suggests}
\label{sec:suggests}

Each is a suggestion rather than a theorem, and i flag which are
checkable.

\paragraph{(i) The import into $\Par$ has a known shape.}
``Behaviour specified in the source category, imported into the par
rule'' is the shape of a distributive law of syntax over behaviour.
Turi and Plotkin's bialgebraic semantics takes operational rules
given as a natural transformation $\lambda : SF \Rightarrow
FS^{\dagger}$, with $S$ a syntax functor, $F$ a behaviour functor
and $S^{\dagger}$ the free monad on $S$, and yields an operational
model together with a canonical, internally fully abstract
denotational model \cite{turiplotkin}.
The relevant corollary is that for any such law, behavioural
equivalence in the final coalgebra is a congruence on the
operational model.
If the import can be put in that form, bisimulation-as-congruence
over $\RHO[\Bee]$ comes free rather than being re-proved per
$\Bee$.
If it cannot, congruence fails and every compositional result of
the Turn stops transferring across the cut.
This is checkable and unchecked, and it is the single most
consequential unchecked thing in this chapter.

\paragraph{(ii) Unambiguity forces behaviours, not states.}
If the objects of $\Bee$ are states with dynamics, then $\quo{b}$
codes a moving target: the same mark has different compliants at
different times, which is exactly the failure of
Condition~\ref{cond:ua}.
The repair is forced.
The objects of $\Bee$ must be \emph{behaviours} --- a subobject of
the final $F$-coalgebra --- so that quoting a builtin quotes its
entire future rather than a snapshot.
Then $*\quo{b} \bisim b$ holds definitionally, and unambiguity
holds because behaviours do not change over time even though states
do.
A consequence worth stating on its own: one cannot take $\Bee$ to
be the physical furniture of the world.
$\Bee$ must be the world already quotiented by an observational
equivalence --- which is where Conjecture~\ref{conj:fd} would
enter, with $\Bee = \Qs(W)$, and which ties the whole construction
back to a budget.

\paragraph{(iii) The source category is constrained.}
Name matching in $\inn{y}{x}{P} \Par \out{x}{Q}$ requires deciding
name equality, and names now include $\quo{b}$.
So structural congruence on $\RHO[\Bee]$ is decidable only if
equality on $\Bee$'s objects is at least budget-decidable.
This is a restriction on which worlds can serve as source, and it
is the formal descendant of the fact that builtin names are the
base case of the well-founded coding by which quotation
manufactures fresh names.

\section{Why this is an analogy and not a solution}
\label{sec:notasolution}

i now argue against the construction just presented, because the
argument is short and decisive and it is better made here than by a
reader.

\begin{enumerate}[label=\textbf{(\alph*)},leftmargin=2.4em]

\item \textbf{It consumes what it was meant to produce.}
$\RHO[-]$ takes $\Bee$ as an argument.
The reference bootstrap problem is the problem of \emph{finding}
$\Bee$: of individuating the world into behaviours worth having
symbols for.
The functor presupposes a solved instance of the problem it was
introduced to address, and everything it then does is bookkeeping
over that solution.
This is the decisive objection and the rest are supporting.
It is also, i note without pleasure, the same objection as
Observation~\ref{obs:sidestep}, one level of abstraction up: the
first time we assumed the world was made of processes, and this
time we assume it comes pre-carved into them.

\item \textbf{Functoriality entails relabelling invariance.}
$\RHO[\Bee] \cong \RHO[\Bee']$ whenever $\Bee \cong \Bee'$.
So the construction cannot fix reference; it transports whatever
reference $\Bee$ came with.
This is not an imported objection but a theorem of the
construction, and it is the formal version of
Remark~\ref{rem:invariance} and the machinery behind
Observation~\ref{obs:social}.

\item \textbf{Losslessness is stipulated where it is interesting.}
$*\quo{b} \bisim b$ holds by construction, because we defined
$\quo{-}$ on $\Bee$ to be exact.
In the world, sensor-to-symbol transduction is lossy, noisy and
partial, and the interesting question is what to do when the
roundtrip \emph{fails}.
The construction has nothing to say about approximate compliance,
which is where all the empirical difficulty lives.

\item \textbf{The par rule import is unverified.}
Section~\ref{sec:suggests}(i) is conditional.
Until the import is exhibited as a distributive law of the required
shape, congruence of bisimulation over $\RHO[\Bee]$ is an
assumption, and with it every compositional result that would be
transported.

\item \textbf{Two interaction disciplines complicate the logic.}
Communication in $\rhoc$ is name-mediated.
If builtins react in the source category by something that is not a
name rendezvous, then $\RHO[\Bee]$'s parallel composition hosts two
kinds of redex, and the namespace logic's modalities --- which
quantify over redex positions --- require a sort distinction.
The ladders computed in Chapters~\ref{ch:gam} and~\ref{ch:com}
would need recomputation, and it is not known whether their
qualitative shape survives.

\item \textbf{Nothing here proposes candidates.}
Even granting all of the above, the construction is a specification
language for grounded symbol systems, not a procedure for acquiring
one.
It says what the answer must look like.
It does not search.

\end{enumerate}

\noindent
The honest summary is that the functorial reading is a good way to
\emph{state} several of Goodman's conditions in a form where they
might be proved, and no way at all to satisfy them.

\section{Where a transformer belongs}
\label{sec:transformer}

The framing above assigns transformers a job, and it is not the job
they are usually given in discussions of grounding.
It is also not quite the job Chapter~\ref{ch:sci} gave them, and
i have corrected that chapter rather than leave the two accounts
standing side by side.

The slogan there was: the network is the transducer, the ecology is
the epistemology.
The second half is right.
The first half is wrong, and it is wrong in an instructive
direction.
In the construction of Section~\ref{sec:functor} the transducer
\emph{is} $\quo{-}$ and $*$ extended to builtins.
That map is exact and lossless by construction; a network cannot
improve on it and is not needed for it.
What is needed, and what nothing in the calculus supplies, is a
\emph{proposal mechanism}: given a stream of host-world
observation, propose candidate builtin behaviours.
That is objection~(a) of Section~\ref{sec:notasolution} turned into
a task, which is the most useful thing one can do with a decisive
objection.

Two features make it a well-posed learning problem, which is more
than can be said for ``learn symbols from sensors''.

\paragraph{The success criterion is coverage, not reward.}
$\Bee$ is a \emph{presentation}: a generating set such that the
observable world lies in the image of $\RHO[\Bee]$.
This is a generators-and-relations problem with a definite answer,
and minimality of $\Bee$ is a statistical-complexity analogue
rather than a hyperparameter to be tuned.

\paragraph{There is a principled candidate.}
Causal states are already the coarsest coarse-graining of pasts
retaining full predictive power \cite{shalizi2001}, and by
Section~\ref{sec:recovers} a next-token-trained transformer
represents belief states over them \cite{shai2024}.
That makes causal states the natural first candidate for which
distinctions deserve builtin names.
The correspondence extends further than convenience: Crutchfield's
taxonomy of $\epsilon$-machines runs from finite-state through
countable to uncountable and fractal, and that taxonomy \emph{is} a
density hierarchy --- so Condition~\ref{cond:sfd} is the condition
separating its bottom rung from the rest, and the question ``is
this world notationally accessible to this agent at this budget''
becomes ``is the budget-quotiented $\epsilon$-machine of the host
process finite''.

That last equivalence is, i think, the most promising single item
in this chapter, because both sides of it are computable and
neither side has been computed.

A second job for a network is amortised bisimulation checking.
Assays are priced in the framework already, checking behavioural
equivalence is expensive, and an approximate checker with
calibrated error is a legitimate purchase if its errors can be
priced.
That is more speculative and i do not pursue it.

\section{Open questions}
\label{sec:notation-open}

These are the ones i would take money on being tractable.

\begin{question}[Transitivity of budget-relative indistinguishability]
\label{q:trans}
Is $\bisim^{\sigma}$ of Conjecture~\ref{conj:fd} transitive?
If not, what is the cost of taking transitive closure, and does the
resulting scheme still separate at a rate related to $\sigma$?
\end{question}

\begin{question}[Density: agent or world?]
\label{q:density}
Given a failure of Condition~\ref{cond:sfd}, is there a criterion
separating ``this agent cannot notate this world'' from ``this
world admits no notation''?
A candidate: the former is a statement about $\Qs$ for the agent's
$\sigma$, the latter about $\lim_{\sigma \to \infty} \Qs$.
\end{question}

\begin{question}[Notationality of learned codebooks]
\label{q:codebooks}
Goodman's five conditions are checkable properties of an
encoder / decoder pair.
What do VQ-VAE codebooks, BPE tokenisers and LatPlan's state
autoencoder score on them?
``How notational is this tokeniser'' appears to be a measurable
quantity and, so far as i can determine, has never been measured.
\end{question}

\begin{question}[Is the import a distributive law?]
\label{q:gsos}
Can the import of $\Bee$'s behaviour into parallel composition be
exhibited as $\lambda : \RHO \circ F \Rightarrow F \circ
\RHO^{\dagger}$?
If so, congruence is free by \cite{turiplotkin}.
If not, what weaker condition holds, and does the lax version
suffice?
\end{question}

\begin{question}[Does $\RHO$ preserve the cost-accounting monad?]
\label{q:monad}
The Turn works inside plain $\rhoc$ and appeals to a WLOG on syntax
that is valid there.
Do $\RHO[-]$ and the cost monad of Chapter~\ref{ch:costmonad}
commute?
If they do not, the foraging inequality, the pricing schedule and
the monotonicity of cost in distance do not transport to
$\RHO[\Bee]$, and with them Conjecture~\ref{conj:fd} loses its
setting.
\end{question}

\begin{question}[The minimal social configuration]
\label{q:social}
Observation~\ref{obs:social} says grounding is a property of a
composite.
What is the smallest composite that suffices?
Two agents with shared action, or does the argument require a
population?
Does the typed-overlap apparatus of Chapter~\ref{ch:com} already
contain the answer?
\end{question}

\begin{question}[Approximate compliance]
\label{q:approx}
Section~\ref{sec:notasolution}(c) observes that the interesting
case is roundtrip failure.
Is there a notion of \emph{graded} compliance that degrades
Goodman's conditions gracefully without collapsing into the metric
of Observation~\ref{obs:metric}?
\end{question}

\section{What is claimed and what is not}
\label{sec:notation-claims}

\paragraph{Claimed.}
That the correspondence between a language and a world is a fact
about the language's users and is not contained in the corpus
(Observation~\ref{obs:locus}); that the record of decipherment and
of cetacean bioacoustics is evidence for this rather than merely
consistent with it, and that the two successful routes both acquire
the correspondence from agents who already have it
(Observation~\ref{obs:transport}); that next-token prediction
recovers the structure of the token-generating process, which for a
corpus is the linguistic community and not the environment
(Observation~\ref{obs:generator}); that the problem has at least
two independent solutions in nature, so that solubility is not in
question and impossibility arguments are unsound
(Observation~\ref{obs:soluble}); that this book's grounding is
analytic and therefore does not address the problem
(Observation~\ref{obs:sidestep}); that Goodman's conditions
decompose into quotient conditions bisimulation supplies and
separation conditions it does not (Observation~\ref{obs:sep}); that
bisimulation metrics are the wrong repair for the latter
(Observation~\ref{obs:metric}); that quotation and dereference in
the pure calculus are score-internal recursion and acquire a
compliance class outside the syntax only when builtin processes are
adjoined (Observations~\ref{obs:tworeadings}
and~\ref{obs:enters}); that Goodman's adequacy criterion is
hosting-and-exhausting for the unit of $\RHO[-]$
(Observation~\ref{obs:hosting}); and that no relabelling-invariant
single-agent framework can deliver correspondence to \emph{our}
symbols (Observation~\ref{obs:social}).

\paragraph{Conjectured.}
That cost accounting supplies finite differentiation with the
budget as separation constant (Conjecture~\ref{conj:fd}).

\paragraph{Offered as analogy only.}
The entire content of
Sections~\ref{sec:rho-builtins}--\ref{sec:functor}.
The functorial reading organises the conditions attractively and
suggests several checkable statements.
It does not bootstrap anything, for the reasons in
Section~\ref{sec:notasolution}, and the first of those reasons is
not a technicality.

\paragraph{Not claimed.}
That Goodman's conditions are the right ones --- they are adopted
provisionally, and the density and wrong-note objections are live.
That transformers cannot in principle transduce --- only that
next-token prediction over an already-grounded corpus is not
transduction, and that the two are routinely conflated.
And that the reference bootstrap problem is soluble \emph{by
construction}: nature's two solutions were achieved by populations,
under selection, over long timescales, with interfaces that
co-evolved with the symbol systems they support, and this book has
no argument that any of those conditions can be dispensed with.
